% Part: incompleteness
% Chapter: arithmetization-syntax
% Section: proofs-in-ax

\documentclass[../../../include/open-logic-section]{subfiles}

\begin{document}

\olfileid[hi]{inc}{art}{pax}
\olsection{अभिगृहीतीय \usetoken{P}{derivation}}

\begin{explain}
अभिगृहीतीय !!{derivation} का अंकगणितीकरण करने के लिए हमें !!{derivation} को
संख्याओं के रूप में निरूपित करना होगा। चूँकि !!{derivation} केवल !!{formula}
के अनुक्रम हैं, स्वाभाविक उपाय प्रत्येक !!{derivation} को उसमें आए !!{formula}
के कूटों के अनुक्रम के कूट के रूप में कूटित करना है।
\end{explain}

\begin{defn}
यदि $\delta$, !!{formula} $!A_1$, \dots,~$!A_n$ से बनी अभिगृहीतीय
!!{derivation} है, तो $\Gn{\delta}$ है
\[
\tuple{\Gn{!A_1}, \dots, \Gn{!A_n}}.
\]
\end{defn}

\begin{ex}
  इस बहुत सरल !!{derivation} पर विचार करें:
  \begin{derivation}
  1. & $!B \lif (!B \lor !A)$ \\
  2. & $(!B \lif (!B \lor !A)) \lif (!A  \lif (!B \lif (!B \lor !A)))$\\
  3. & $!A  \lif (!B \lif (!B \lor !A))$
  \end{derivation}
  इस !!{derivation} की ग्योडेल (G\"odel) संख्या होगी
  \begin{align*}
  \openTuple\, 
    & \Gn{!B \lif (!B \lor !A)}, \\
    &\Gn{(!B \lif (!B \lor !A)) \lif (!A  \lif (!B \lif (!B \lor !A)))},\\
    & \Gn{!A  \lif (!B \lif (!B \lor !A))} \,\closeTuple.
  \end{align*}
\end{ex}

\begin{explain}
!!{derivation} का निरूपण तय कर लेने पर हमें यह भी दिखाना होगा कि ऐसी
!!{derivation} पर आदिम पुनरावर्ती ढंग से क्रियाएँ की जा सकती हैं और उनके
आवश्यक गुणों तथा संबंधों को भी उसी ढंग से व्यक्त किया जा सकता है। कुछ
संक्रियाएँ सरल हैं: उदाहरणतः किसी !!a{derivation} की ग्योडेल (G\"odel)
संख्या~$d$ दिए जाने पर $(d)_{\len{d}-1}$ उसके अंतिम-!!{formula} की ग्योडेल (G\"odel)
 संख्या देता है। कुछ संक्रियाएँ कहीं अधिक कठिन हैं। हम कम-से-कम यह
रूपरेखा देंगे कि उन्हें कैसे किया जाए। लक्ष्य यह दिखाना है कि ``$\delta$,
~$\Gamma$ से~$!A$ की !!a{derivation} है'' संबंध, $\delta$ और~$!A$ की ग्योडेल (G\"odel)
 संख्याओं में आदिम पुनरावर्ती है।
\end{explain}

\begin{prop}
\ollabel{prop:followsby}
निम्न संबंध आदिम पुनरावर्ती हैं:
\begin{enumerate}
\item $!A$ एक अभिगृहीत है।
\item $\delta$ की $i$-वीं पंक्ति मोडस पोनेन्स से औचित्यप्राप्त है।
\item $\delta$ की $i$-वीं पंक्ति \QR{} से औचित्यप्राप्त है।
\item $\delta$ एक सही !!{derivation} है।
\end{enumerate}
\end{prop}

\begin{proof}
हमें दिखाना है कि !!{formula} की ग्योडेल (G\"odel) संख्याओं और !!{derivation}
की ग्योडेल (G\"odel) संख्याओं के बीच संगत संबंध आदिम पुनरावर्ती हैं।
\begin{enumerate}
\item हमारे पास अभिगृहीत-प्रतिरूपों की एक दी हुई सूची है, और यदि $!A$ इनमें
  से किसी प्रतिरूप द्वारा दिए गए रूप का है, तो वह अभिगृहीत है। चूँकि
  प्रतिरूपों की सूची परिमित है, इतना दिखाना पर्याप्त है कि प्रत्येक
  अभिगृहीत-प्रतिरूप के लिए हम आदिम पुनरावर्ती ढंग से जाँच सकते हैं कि $!A$ उस
  रूप का है या नहीं। उदाहरण के लिए, अभिगृहीत-प्रतिरूप पर विचार करें
  \[
  !B \lif (!C \lif !B).
  \]
  $!A$ इस अभिगृहीत-प्रतिरूप का उदाहरण है यदि ऐसे !!{formula} $!B$ और $!C$
  हैं कि `$($', $!B$, `$\lif$', `$($', $!C$, `$\lif$', $!B$ और~`$))$' को
  जोड़ने पर $!A$ मिलता है। हम~$!A$ की ग्योडेल (G\"odel) संख्या $n$ का संगत
  गुण जाँच सकते हैं, क्योंकि अनुक्रमों का संयोजन आदिम पुनरावर्ती है और $!B$
  तथा $!C$ की ग्योडेल (G\"odel) संख्याएँ~$!A$ की ग्योडेल (G\"odel) संख्या
  से छोटी होनी चाहिए; संबंध लागू होने पर $!B$ और $!C$, दोनों~$!A$ के
  उप-!!{formula} हैं। अतः हम परिभाषित कर सकते हैं:
  \begin{multline*}
  \fn{IsAx}_{!B \lif (!C \lif !B)}(n) \defiff \bexists{b< n}{
    \bexists{c < n}{(\fn{Sent}(b) \land \fn{Sent}(c) \land {}}}\\ n =
  \Gn{(} \concat b \concat \Gn{\lif} \concat \Gn{(} \concat c \concat
  \Gn{\lif} \concat b \concat \Gn{))}).
  \end{multline*}
  यदि प्रत्येक अभिगृहीत-प्रतिरूप के लिए ऐसी परिभाषा हो, तो उनका वियोजन गुण
  $\fn{IsAx}(n)$, ``$n$ किसी अभिगृहीत की ग्योडेल (G\"odel) संख्या है'',
  परिभाषित करता है।
\item $\delta$ की $i$-वीं पंक्ति मोडस पोनेन्स से औचित्यप्राप्त है यदि और
  केवल यदि ऐसी पंक्तियाँ $j$ और $k < i$ हैं जिनमें पंक्ति~$j$ का !!{sentence}
  कोई सूत्र $!A$, पंक्ति~$k$ का वाक्य $!A \lif !B$, और पंक्ति~$i$ का वाक्य
  ~$!B$ है।
  \begin{multline*}
    \fn{MP}(d, i) \defiff \bexists{j < i}{\bexists{k < i}{}}\\
    (d)_k = \Gn{(} \concat (d)_j
      \concat \Gn{\lif} \concat (d)_i \concat \Gn{)}
  \end{multline*}
  चूँकि परिबद्ध परिमाणीकरण, संयोजन और $=$ आदिम पुनरावर्ती हैं, यह एक आदिम
  पुनरावर्ती संबंध परिभाषित करता है।
\item $\delta$ की कोई पंक्ति \QR{} से औचित्यप्राप्त है यदि वह $!B \lif
  \lforall[x][!A(x)]$ रूप की है, कोई पिछली पंक्ति किसी !!{constant}~$c$ के
  लिए $!B \lif !A(c)$ है, और $c$,~$!B$ में नहीं आता। यह लागू होता है यदि और
  केवल यदि
  \begin{enumerate}
  \item कोई !!a{sentence}~$!B$ है, और
  \item केवल एक चर~$x$ मुक्त रखने वाला !!a{formula}~$!A(x)$ ऐसा है कि
  \item पंक्ति $i$ में $!B \lif \lforall[x][!A(x)]$ है,
  \item किसी पंक्ति $j < i$ में किसी अचर~$c$ के लिए $!B \lif
    \Subst{!A}{c}{x}$ है,
  \item और वह अचर~$!B$ में नहीं आता।
  \end{enumerate}
  इन सभी को आदिम पुनरावर्ती ढंग से जाँचा जा सकता है, क्योंकि $!B$, $!A(x)$
  और $x$ की ग्योडेल (G\"odel) संख्याएँ पंक्ति~$i$ के सूत्र की ग्योडेल (G\"odel)
   संख्या से छोटी हैं, और $a$ की संख्या पंक्ति~$j$ के सूत्र की
  ग्योडेल (G\"odel) संख्या से छोटी है:
  \begin{multline*}
    \fn{QR}_1(d, i) \defiff \bexists{b < (d)_i}{\bexists{x <
        (d)_i}{\bexists{a < (d)_i}{\bexists{c < (d)_j}{(}}}}
    \\ \fn{Var}(x) \land \fn{Const}(c) \land {} \\
    (d)_i = \Gn{(} \concat
    b \concat \Gn{\lif} \concat \Gn{\lforall} \concat x \concat a
    \concat \Gn{)} \land {}\\
    (d)_j = \Gn{(} \concat b \concat
    \Gn{\lif} \concat \fn{Subst}(a,c,x) \concat \Gn{)} \land {}\\ \fn{Sent}(b)
    \land \fn{Sent}(\fn{Subst}(a,c,x)) \land {} \bforall{k <
      \len{b}}{(b)_k \neq (c)_0})
  \end{multline*}
  यहाँ हम मानते हैं कि $c$ और $x$, पदों के रूप में लिए गए चर और अचर की
  ग्योडेल (G\"odel) संख्याएँ हैं (अर्थात् उनके चिह्न-कूट नहीं)। हम यह जाँचकर
  परीक्षण करते हैं कि $x$,~$!A(x)$ का एकमात्र मुक्त चर है कि
  $\Subst{!A(x)}{c}{x}$ !!a{sentence} है या नहीं; और~$!B$ के प्रत्येक चिह्न
  को~$c$ से भिन्न अपेक्षित करके सुनिश्चित करते हैं कि $c$,~$!B$ में नहीं आता।

  \QR{} का दूसरा रूप अभ्यास के लिए छोड़ते हैं।
\item $d$ किसी सही !!{derivation} की ग्योडेल (G\"odel) संख्या है यदि और केवल
  यदि उसकी प्रत्येक पंक्ति अभिगृहीत है अथवा मोडस पोनेन्स या~\QR{} से
  औचित्यप्राप्त है। अतः:
  \[
  \fn{Deriv}(d) \defiff \bforall{i < \len{d}}{(\fn{IsAx}((d)_i) \lor
    \fn{MP}(d,i) \lor \fn{QR}(d, i))}
  \]
\end{enumerate}
\end{proof}

\begin{prob}
\olref[inc][art][pax]{prop:followsby} की भाँति निम्न संबंध परिभाषित करें:
\begin{enumerate}
\item $\fn{IsAx}_{!A \lif (!B \lif (!A \land !B))}(n)$,
\item $\fn{IsAx}_{\lforall[x][!A(x)] \lif !A(t)}(n)$,
\item $\fn{QR}_{2}(d, i)$ (\QR{} के दूसरे रूप के लिए)।
\end{enumerate}
\end{prob}

\begin{prop}
मान लें कि $\Gamma$, !!{sentence} का आदिम पुनरावर्ती समुच्चय है। तब यह
व्यक्त करने वाला संबंध $\Prf[\Gamma](x, y)$ कि ``$x$,~$\Gamma$ से~$!A$ की
!!a{derivation}~$\delta$ का कूट है और $y$,~$!A$ की ग्योडेल (G\"odel) संख्या
है'', आदिम पुनरावर्ती है।
\end{prop}

\begin{proof}
मान लें कि ``$y \in \Gamma$'' आदिम पुनरावर्ती विधेय~$R_\Gamma(y)$ से दिया
गया है। हमें दिखाना है कि संबंध $\Prf[\Gamma](x, y)$ आदिम पुनरावर्ती है,
जहाँ $\Prf[\Gamma](x, y)$ तब लागू होता है जब $y$ किसी !!a{sentence}~$!A$ की
ग्योडेल (G\"odel) संख्या और $x$,~$\Gamma$ से~$!A$ की !!a{derivation} का कूट
हो।

पिछली प्रतिज्ञप्ति के अनुसार गुण $\fn{Deriv}(x)$, जो तभी लागू होता है जब $x$
किसी सही !!{derivation}~$\delta$ का कूट हो, आदिम पुनरावर्ती है। किंतु उस
परिभाषा ने !!{derivation} की पंक्तियों को औचित्य देने के अतिरिक्त उपाय के रूप
में समुच्चय $\Gamma$ को सम्मिलित नहीं किया। किसी पंक्ति के \QR{} से
औचित्यप्राप्त होने के हमारे आदिम पुनरावर्ती परीक्षण ने यह अपेक्षा भी छोड़ दी
कि अचर~$c$ को~$\Gamma$ में आने की अनुमति नहीं। परिभाषा में संशोधन करके
$\Gamma$ को सीधे सम्मिलित करना संभव है, किंतु $\fn{Deriv}$ और निगमन प्रमेय का
उपयोग अधिक सरल है। $\Gamma \Proves !A$ यदि और केवल यदि !!{sentence} की कोई
परिमित सूची $!B_1$, \dots, $!B_n \in \Gamma$ ऐसी है कि $\{!B_1, \dots,
!B_n\} \Proves !A$। और निगमन प्रमेय के अनुसार यह तब लागू होता है जब $\Proves
(!B_1 \lif (!B_2 \lif \cdots (!B_n \lif !A)\cdots))$। ग्योडेल (G\"odel)
संख्या~$z$ वाला कोई !!a{sentence} इस रूप का है या नहीं, इसे आदिम पुनरावर्ती
ढंग से जाँचा जा सकता है। इसलिए $x$ को~$\Gamma$ \emph{से} ग्योडेल (G\"odel)
संख्या~$y$ वाले !!{sentence} की !!a{derivation} की ग्योडेल (G\"odel) संख्या
मानने के स्थान पर हम $x$ को~$\emptyset$ से उपर्युक्त रूप के अंतर्निहित
प्रतिबंधात्मक की !!a{derivation} की ग्योडेल (G\"odel) संख्या मानते हैं।

पहले, यदि हमारे पास !!{sentence} का कोई अनुक्रम है, तो हम आदिम पुनरावर्ती
ढंग से ऐसा प्रतिबंधात्मक बना सकते हैं जिसमें ये सभी वाक्य पूर्वपक्ष और दिया
हुआ !!{sentence} उत्तरपक्ष हो:
\begin{align*}
  \fn{hCond}(s, y, 0) & = y \\
  \fn{hCond}(s, y, n+1) & = \Gn{(} \concat (s)_{n} \concat \Gn{\lif}
  \concat \fn{Cond}(s, y, n) \concat \Gn{)}\\
  \fn{Cond}(s, y) & = \fn{hCond}(s, y, \len{s})\\
  \intertext{अतः हम $\Prf[\Gamma](x, y)$ को इस प्रकार परिभाषित कर सकते हैं:}
  \Prf[\Gamma](x, y) & \defiff \bexists{s < \fn{sequenceBound}(x,x)}{(} \\
  &\qquad (x)_{\len{x}-1} = \fn{Cond}(s,y) \land {} \\
  &\qquad\bforall{i<\len{s}}{(s)_i \in \Gamma} \land {} \\
  &\qquad\fn{Deriv}(x)).
\end{align*}
~$s$ पर परिसीमा इस बात से मिलती है कि प्रत्येक $(s)_i$, !!{derivation} की
अंतिम पंक्ति के किसी उप-!!{formula} की ग्योडेल (G\"odel) संख्या है; अर्थात्
वह $(x)_{\len{x}-1}$ से छोटा है। पूर्वपक्षों $!B \in \Gamma$ की संख्या,
अर्थात्~$s$ की लंबाई,~$x$ की अंतिम पंक्ति की लंबाई से छोटी है।
\end{proof}

\end{document}
